\section{Entorno de pruebas}

La validación del sistema se ha llevado a cabo en tres escenarios de despliegue
de complejidad creciente, que reproducen casos de uso reales de producción audiovisual
en directo.

\subsection{Escenario 1 --- Monopuesto}

En este escenario, voctocore, las fuentes FFmpeg y voctogui se ejecutan en un único
equipo conectado en loopback. Es el más sencillo de desplegar y resulta adecuado para
desarrollo, depuración y demostraciones. El guion de automatización
\texttt{1pc\_escenario1/pruebas\_lanzar\_realizacion.sh} encadena el arranque de todos
los componentes en el orden correcto: primero voctocore, luego las fuentes (con espera
activa del puerto 9999) y finalmente la interfaz gráfica.

El principal factor limitante es el ancho de banda interno: una señal UYVY a
1920×1080 @25\,fps genera aproximadamente 830\,Mbps por fuente sobre la interfaz de
\textit{loopback} TCP, lo que requiere una CPU de al menos cuatro núcleos a 3\,GHz
para gestionar cuatro fuentes simultáneas.

\subsection{Escenario 2 --- Cliente-servidor en red local}

En este escenario, voctocore y las fuentes se ejecutan en un PC dedicado
(\textit{servidor/cerebro}), mientras que voctogui se conecta desde un segundo
equipo mediante la red local. El guion \texttt{arrancar\_voctocore\_pc1.sh} detecta
automáticamente la dirección IP local del servidor y la expone para que el realizador
la introduzca en su equipo. La GUI del realizador (\texttt{arrancar\_voctogui\_pc2.sh})
se conecta al servidor por TCP y recibe el stream de previsualización comprimido en
JPEG, minimizando el ancho de banda requerido entre ambos equipos.

Este escenario permite separar físicamente el rack de procesamiento de la mesa de
realización, distribuyendo la carga de CPU entre dos equipos.

\subsection{Escenario 3 --- Docker Compose}

El despliegue completo mediante Docker Compose se realiza con un único comando
(\texttt{docker compose up -d}) desde el directorio \texttt{docker\_escenario3/}.
Docker Compose levanta los diez servicios definidos en \texttt{docker-compose.yml}
en el orden correcto, gracias a las dependencias (\texttt{depends\_on}) y los
\textit{healthchecks} configurados. El realizador se conecta desde un segundo equipo
exactamente igual que en el Escenario 2.

Este escenario es el recomendado para despliegues en producción real por su máxima
reproducibilidad: la imagen Docker garantiza que el entorno de ejecución es idéntico
independientemente del sistema operativo anfitrión.

% --- Figura sugerida ---
% \begin{figure}[H]
%   \centering
%   \includegraphics[width=0.9\textwidth]{figuras/diagramas_escenarios.png}
%   \caption{Diagramas de red de los tres escenarios de despliegue.}
%   \label{fig:diagramas_escenarios}
% \end{figure}

% --- Tabla sugerida ---
% \begin{table}[H]
%   \centering
%   \caption{Especificaciones hardware de los equipos empleados en las pruebas.}
%   \label{tab:hardware_pruebas}
%   \begin{tabular}{|l|l|l|}
%     \hline
%     \textbf{Equipo} & \textbf{CPU} & \textbf{RAM / Red} \\
%     \hline
%     PC Principal (servidor)  & Intel Core i7-xxxx @ x núcleos & xx GB / 1 GbE \\
%     PC Realizador (cliente)  & Intel Core i5-xxxx @ x núcleos & xx GB / 1 GbE \\
%     \hline
%   \end{tabular}
% \end{table}
